% Exam Template for UMTYMP and Math Department courses
%
% Using Philip Hirschhorn's exam.cls: http://www-math.mit.edu/~psh/#ExamCls
%
% run pdflatex on a finished exam at least three times to do the grading table on front page.
%
%%%%%%%%%%%%%%%%%%%%%%%%%%%%%%%%%%%%%%%%%%%%%%%%%%%%%%%%%%%%%%%%%%%%%%%%%%%%%%%%%%%%%%%%%%%%%%

% These lines can probably stay unchanged, although you can remove the last
% two packages if you're not making pictures with tikz.
\documentclass[11pt]{exam}
\RequirePackage{amssymb, amsfonts, amsmath, latexsym, verbatim, xspace, setspace}
\RequirePackage{tikz, pgflibraryplotmarks}

% By default LaTeX uses large margins.  This doesn't work well on exams; problems
% end up in the "middle" of the page, reducing the amount of space for students
% to work on them.
\usepackage[margin=1in]{geometry}
\usepackage{enumerate}
\usepackage{amsthm}

\theoremstyle{definition}
\newtheorem{soln}{Solution}

% Here's where you edit the Class, Exam, Date, etc.
\newcommand{\class}{Math 407 Section 1}
\newcommand{\term}{Fall 2025}
\newcommand{\examnum}{Final Exam}
\newcommand{\examdate}{May 14, 2025}
\newcommand{\timelimit}{1 Hour 50 Minutes}
\newcommand{\ol}[1]{\overline{#1}}
\newcommand{\wt}[1]{\widetilde{#1}}

% For an exam, single spacing is most appropriate
\singlespacing
% \onehalfspacing
% \doublespacing

% For an exam, we generally want to turn off paragraph indentation
\parindent 0ex

\begin{document} 

% These commands set up the running header on the top of the exam pages
\pagestyle{head}
\firstpageheader{}{}{}
\runningheader{\class}{\examnum\ - Page \thepage\ of \numpages}{\examdate}
\runningheadrule

\begin{flushright}
\begin{tabular}{p{2.8in} r l}
\textbf{\class} & \textbf{Name (Print):} & \makebox[2in]{\hrulefill}\\
\textbf{\term} &&\\
\textbf{\examnum} & \textbf{Student ID:}&\makebox[2in]{\hrulefill}\\
\textbf{\examdate} &&\\
\textbf{Time Limit: \timelimit} % & Teaching Assistant & \makebox[2in]{\hrulefill}
\end{tabular}\\
\end{flushright}
\rule[1ex]{\textwidth}{.1pt}


This exam contains \numpages\ pages (including this cover page) and
\numquestions\ problems.  Check to see if any pages are missing.  Enter
all requested information on the top of this page, and put your initials
on the top of every page, in case the pages become separated.\\

You may \textit{not} use your books or notes on this exam.\\

You are required to show your work on each problem on this exam.  The following rules apply:\\

\begin{minipage}[t]{3.7in}
\vspace{0pt}
\begin{itemize}

%\item \textbf{If you use a ``fundamental theorem'' you must indicate this} and explain
%why the theorem may be applied.

\item \textbf{Organize your work}, in a reasonably neat and coherent way, in
the space provided. Work scattered all over the page without a clear ordering will 
receive very little credit.  

\item \textbf{Mysterious or unsupported answers will not receive full
credit}.  A correct answer, unsupported by calculations, explanation,
or algebraic work will receive no credit; an incorrect answer supported
by substantially correct calculations and explanations might still receive
partial credit.  This especially applies to limit calculations.

\item If you need more space, use the back of the pages; clearly indicate when you have done this.

\item If the problem asks for a proof, be sure to carefully justify your work, including any theorems from class.

Note: you may NOT use a theorem or result from class to prove something when it makes the problem entirely trivial.  If you are unsure whether a particular theorem or result is allowed, just ask!

Unless otherwise stated, all rings will be commutative with identity and all ring homomorphisms will send the identity to the identity.

\end{itemize}

Do not write in the table to the right.
\end{minipage}
\hfill
\begin{minipage}[t]{2.3in}
\vspace{0pt}
%\cellwidth{3em}
\gradetablestretch{2}
\vqword{Problem}
\addpoints % required here by exam.cls, even though questions haven't started yet.	
\gradetable[v]%[pages]  % Use [pages] to have grading table by page instead of question

\end{minipage}
\newpage % End of cover page

%%%%%%%%%%%%%%%%%%%%%%%%%%%%%%%%%%%%%%%%%%%%%%%%%%%%%%%%%%%%%%%%%%%%%%%%%%%%%%%%%%%%%
%
% See http://www-math.mit.edu/~psh/#ExamCls for full documentation, but the questions
% below give an idea of how to write questions [with parts] and have the points
% tracked automatically on the cover page.
%
%
%%%%%%%%%%%%%%%%%%%%%%%%%%%%%%%%%%%%%%%%%%%%%%%%%%%%%%%%%%%%%%%%%%%%%%%%%%%%%%%%%%%%%

\begin{questions}

\addpoints

\question[10]\mbox{}
\textbf{TRUE or FALSE!}  Write  TRUE if the statement is true.  Otherwise, write FALSE.  Your response should be in ALL CAPS.  No justification is required.
\begin{enumerate}[(a)]
\item  
The discriminant of a cubic polynomial can tell us if its roots are all real
\vspace{1.3in}
\item  
If $F\subseteq K$ is a field extension, then $\text{Gal}_F(K) = [K:F]$
\vspace{1.3in}
\item  
$\mathbb Q(\sqrt[3]{5})$ is a splitting field
\vspace{1.3in}
\item  
The ring $\mathbb Z[\sqrt{-5}]$ is a UFD
\vspace{1.3in}
\item  
Any subgroup of a cyclic group is also cyclic
\vspace{1.3in}
\end{enumerate}

\newpage
\question[10]\mbox{}

\begin{enumerate}[(a)]
\item State the definition of a principal domain
\vspace{3in}
\item State the third isomorphism theorem for groups
\vspace{3in}
\item Write down two non-isomorphic non-Abelian groups of order $8$
\end{enumerate}

\newpage
\question[10]\mbox{}
\begin{enumerate}[(a)]
\item  Give an example of a field extension of $\mathbb Q$ of degree $2$
\vspace{2in}
\item  Give an example of a field extension of $\mathbb Q$ which is transcendental
\vspace{2in}
\item  Give an example of a field extension of $\mathbb Q$ which is not normal
\vspace{2in}
\item  Give an example of a field extension of $\mathbb Q$ whose Galois group is isomorphic to $S_3$
\end{enumerate}

\newpage
\question[10]\mbox{}

Consider the group
$$G = \mathbb Z_{12}\times\mathbb Z_{12}.$$

\begin{enumerate}[(a)]
\item Determine all the elements in the subgroup $H$ generated by $(1,1)$ and $(0,3)$.
\vspace{3in}
\item Explicitly describe the cosets of $G/H$.  How many are there?
\vspace{3in}
\item What standard group is $G/H$ isomorphic to?
\end{enumerate}


\newpage
\question[10]\mbox{}

\begin{enumerate}[(a)]
\item State the Structure Theorem for Finite Abelian Groups
\vspace{3in}
\item Write the group $\mathbb Z_{35}\times \mathbb Z_{15}\times\mathbb Z_{100}$ in invariant factor form
\vspace{3in}
\item how many Abelian groups of order $2^6$ are there, up to isomorphism?
\end{enumerate}

\newpage
\question[10]\mbox{}

Consider the set

$$R = \left\lbrace\frac{a}{b}: a,b\in\mathbb Z,\ b\ \text{odd}\right\rbrace\subseteq \mathbb Q.$$

\begin{enumerate}[(a)]
\item Prove that the set $R$ is a subring of $\mathbb Q$.
\vspace{3in}
\item Prove that $\frac{a}{b}\in R$ is a unit if and only if $a$ is odd.
\vspace{3in}
\item Prove that the principal ideal $(2)\subseteq R$ is a maximal ideal.
\end{enumerate}

\newpage
\question[10]\mbox{}

Let $\mathbb Q\subseteq K\subseteq \mathbb C$ be a splitting field of $(x^2-2)(x^2-3)$.

\begin{enumerate}[(a)]
\item Prove that $K = \mathbb Q(\sqrt{2}, \sqrt{3})$ and $[K:\mathbb Q] = 4$
\vspace{3in}
\item Explain why the Galois group $\text{Gal}_{\mathbb Q}(K)$ has four elements
\vspace{1in}
\item Prove that if $\sigma\in \text{Gal}_{\mathbb Q}(K)$ then $\sigma(\sqrt{2}) = \pm \sqrt{2}$ and $\sigma(\sqrt{3}) = \pm \sqrt{3}$.
\vspace{2in}
\item Prove that $\text{Gal}_{\mathbb Q}(K)$ is isomorphic to $\mathbb Z_2\times\mathbb Z_2$
\end{enumerate}







\end{questions}
\end{document}
